\begin{frame}[fragile]{MDRangePolicy}

  {\Huge Tightly Nested Loops with MDRangePolicy}

  \vspace{20pt}

  \textbf{Learning objectives:}
  \begin{itemize}
    \item{Demonstrate usage of the MDRangePolicy with tightly nested loops.}
    \item{Syntax - Required and optional settings}
    \item{Code demo and example}
  \end{itemize}

  \vspace{-20pt}

\end{frame}

%==========================================================================

\begin{frame}[fragile]{MDRangePolicy (0)}

  \textbf{Motivating example}: Consider the nested for loops:

%  \vspace{5pt}

  \lstset{mathescape, escapeinside={<@}{@>},
          language=C,
          keywords={}}

  %\begin{lstlisting}%[linebackgroundcolor={
    %    \btLstHL{1-3}{darkred!20}
    %  }
    %]

  %for ( int i = 0; i < N0; ++i )
  %for ( int j = 0; j < N1; ++j )
  %for ( int k = 0; k < N2; ++k )
  %  some_init_fcn(i, j, k);

  %\end{lstlisting}

  \begin{code}[linebackgroundcolor={
        \btLstHL<1->{5}{bodyColor}
      },
      frame=single
    ]

@patternfor@pattern (int i = @policy0; i < N0@policy; ++i)
  @patternfor@pattern (int j = @policy0; j < N1@policy; ++j)
    @patternfor@pattern (int k = @policy0; k < N2@policy; ++k)
      some_init_fcn(i, j, k);

  \end{code}

%  \vspace{-11pt}

%  \begin{itemize}
%    \item{Based on Kokkos lessons thus far, you might parallelize this as follows:}
%  \end{itemize}

  {Based on Kokkos lessons thus far, you might parallelize this as}

  \pause

  %\begin{lstlisting}%[linebackgroundcolor={
    %    \btLstHL{3-4}{bodyColor}
    %  }
    %]

  \begin{code}[linebackgroundcolor={
        \btLstHL<1->{5}{bodyColor}
      },
      frame=single
    ]
@patternKokkos::parallel_for@pattern("NestedLoop", @policyN0@policy,
                     KOKKOS_LAMBDA (const i) {
                       @patternfor@pattern (int j = @policy0; j < N1@policy; ++j)
                         @patternfor@pattern (int k = @policy0; k < N2@policy; ++k)
                           some_init_fcn(i, j, k);
                     }
                    );
  \end{code}
  %\end{lstlisting}

%  \begin{textblock*}{0.5\textwidth}(0.08\textwidth,0.235\textheight)
%    \rotatebox{90}{{\footnotesize {\color{darkred!80} section 1}}}
%  \end{textblock*}

%  \begin{textblock*}{0.5\textwidth}(0.08\textwidth,0.40\textheight)
%    \rotatebox{90}{{\footnotesize {\color{blue!80} section 2}}}
%  \end{textblock*}

%  \pause

  \vspace{-5pt}

  \begin{itemize}
    \item{\small{This only parallelizes along one dimension, leaving potential parallelism unexploited.}}
    \item{\small{What if N1 is too small to amortize the cost of constructing a parallel region, but N1*N2*N3 makes it worthwhile?}}
  \end{itemize}

%  \begin{itemize}
%    \item{Where will {\color{darkred!80}section 1} be run?  CPU?  GPU?}
%    \item{Where will {\color{blue!80}section 2} be run?  CPU?  GPU?}
%    \item{How do I \textbf{control} where code is executed?}
%  \end{itemize}

%  \pause
%  \vspace{5pt}

%  \hspace{20pt}{\Large $\Rightarrow$ \textbf{Execution spaces}}

\end{frame}

%==========================================================================

\begin{frame}[fragile]{MDRangePolicy (1)}
   \textbf{OpenMP has a solution: the collapse clause}

  \begin{code}[linebackgroundcolor={
        \btLstHL<1->{5}{bodyColor}
      },
      frame=single
    ]
#pragma @policyomp parallel@policy @patternfor@pattern @policycollapse(3)@policy
@patternfor@pattern (int i = @policy0; i < N0@policy; ++i) {
  @patternfor@pattern (int j = @policy0; j < N1@policy; ++j) {
    @patternfor@pattern (int k = @policy0; k < N2@policy; ++k) {
      /* loop body */
    }
  }
}
  \end{code}

  \pause

Note this changed the policy by adding a `collapse` clause.

\pause

	\vspace{0.5cm}

	\textbf{With Kokkos you also change the policy:}

  \begin{code}[linebackgroundcolor={
        \btLstHL<1->{3}{bodyColor}
      },
      frame=single
    ]
@patternparallel_for@pattern("MDRP", @policyMDRangePolicy<Rank<3>>({0,0,0},{N0,N1,N2})@policy, 
   KOKKOS_LAMBDA(int i, int j, int k) {
     /* loop body */
});
  \end{code}

\end{frame}

\begin{frame}[fragile]{MDRangePolicy (2)}

\begin{block}{MDRangePolicy}
   MDRangePolicy can parallelize tightly nested loops of 1 to 6 dimensions.
\end{block}

	\vspace{0.5cm}

  \begin{code}[linebackgroundcolor={
        \btLstHL<1->{3}{bodyColor}
      },
      frame=single
    ]
@patternparallel_for@pattern("MDRP", @policyMDRangePolicy<Rank<3>>({0,0,0},{N0,N1,N2})@policy, 
   KOKKOS_LAMBDA(int i, int j, int k) {
     /* loop body */
});
  \end{code}

\begin{itemize}
   \item<2-> Specify the dimensionality of the loop with $Rank<DIM>$.
   \item<3-> As with Kokkos Views: only rectangular iteration spaces.
   \item<4-> Provide initializer lists for begin and end values.
   \item<5-> The functor/lambda takes matching number of indicies. 
\end{itemize}
\end{frame}

\begin{frame}[fragile]{MDRangePolicy (3)}
   \textbf{You can also do Reductions:}

  \begin{code}[linebackgroundcolor={
        \btLstHL<1->{5-6}{bodyColor}
      },
      frame=single
    ]
double result;
@patternparallel_reduce@pattern("Label", 
  @policyMDRangePolicy<Rank<3>>({0,0,0},{N0,N1,N2})@policy, 
  KOKKOS_LAMBDA(int i, int j, int k, double& lsum) {
     /* loop body */
  lsum += something;
}, result);
  \end{code}

\begin{itemize}
   \item<2-> The Policy doesn't change the rules for `parallel\_reduce`.
   \item<3-> Additional thread local argument.
   \item<4-> Can do reductions other than "sum" with reducers.
   \item<5-> Can use `View`s as reduction argument.
\end{itemize}
\end{frame}

\begin{frame}[fragile]{MDRangePolicy (4)}
   In structured grid applications a \textbf{tiling} strategy is often used to help with caching.

	\begin{block}{Tiling}
		MDRangePolicy uses a tiling strategy for the iteration space.
	\end{block}

	\begin{itemize}
		\item Specified as a third initializer list.
		\item For GPUs a tile is handled by a single thread block.
			\begin{itemize}
				\item If you provide too large a tile size this will fail!
			\end{itemize}
		\item Optimal tile size is problem and hardware dependent. Tools like Apex could be used to tune the tile size.
	\end{itemize}

\begin{code}[keywords={}]
double result;
parallel_reduce("Label", 
  MDRangePolicy<Rank<3>>({0,0,0},{N0,N1,N2},@policy{T0,T1,T2}@policy), 
  KOKKOS_LAMBDA(int i, int j, int k, double& lsum) {
     /* loop body */
  lsum += something;
}, result);
  \end{code}

\end{frame}


\begin{frame}[fragile]{MDRangePolicy (5)}

  Initializing a Matrix:

\begin{code}[keywords={LayoutLeft}]
View<double**,LayoutLeft> @orangeA@orange("A",@redN0@red,@darkgreenN1@darkgreen);
parallel_for("TestLayoutLeft",
  MDRangePolicy<Rank<2>>({@red0@red,@darkgreen0@darkgreen},{@redN0@red,@darkgreenN1@darkgreen}),
  KOKKOS_LAMBDA(int @redi@red, int @darkgreenj@darkgreen) {
    @orangeA@orange(@redi@red,@darkgreenj@darkgreen) = 1000.0 * @redi@red + 1.0*@darkgreenj@darkgreen;
});
\end{code}

\begin{code}[keywords={LayoutRight}]
View<double**,LayoutRight> @orangeB@orange("B",@redN0@red,@darkgreenN1@darkgreen);
parallel_for("TestLayoutRight",
  MDRangePolicy<Rank<2>>({@red0@red,@darkgreen0@darkgreen},{@redN0@red,@darkgreenN1@darkgreen}),
  KOKKOS_LAMBDA(int @redi@red, int @darkgreenj@darkgreen) {
    @orangeB@orange(@redi@red,@darkgreenj@darkgreen) = 1000.0 * @redi@red + 1.0*@darkgreenj@darkgreen;
});
\end{code}

\pause

	\textbf{How do I make sure that I get the right access pattern?}

\end{frame}


\begin{frame}[fragile]{MDRangePolicy (6)}
\begin{block}{Iteration Pattern}
MDRangePolicy provides compile time control over iteration patterns.
\end{block}

  \begin{lstlisting}[basicstyle=\large,gobble=4]
    Kokkos::Rank< N, IterateOuter, IterateInner >
  \end{lstlisting}

  \begin{itemize}
    \item{\small{\textbf{N: (Required)} the rank of the index space (limited from 1 to 6)}}
    \item{\small{\textbf{IterateOuter (Optional)} iteration pattern between tiles}}
    \begin{itemize}
      \item{\small{\textbf{Options:} Iterate::Left, Iterate::Right, Iterate::Default}}
    \end{itemize}
    \item{\small{\textbf{IterateInner (Optional)} iteration pattern within tiles}}
    \begin{itemize}
      \item{\small{\textbf{Options:} Iterate::Left, Iterate::Right, Iterate::Default}}
    \end{itemize}
  \end{itemize}
\end{frame}

%==========================================================================

\begin{frame}[fragile]{MDRangePolicy (7)}

  Initializing a Matrix fast:

\begin{code}[keywords={LayoutLeft,Iterate,Left}]
View<double**,LayoutLeft> @orangeA@orange("A",@redN0@red,@darkgreenN1@darkgreen);
parallel_for("TestLayoutLeft",
  MDRangePolicy<Rank<2,Iterate::Left,Iterate::Left>>(
	{@red0@red,@darkgreen0@darkgreen},{@redN0@red,@darkgreenN1@darkgreen}),
  KOKKOS_LAMBDA(int @redi@red, int @darkgreenj@darkgreen) {
    @orangeA@orange(@redi@red,@darkgreenj@darkgreen) = 1000.0 * @redi@red + 1.0*@darkgreenj@darkgreen;
});
\end{code}

\begin{code}[keywords={LayoutRight,Iterate,Right}]
View<double**,LayoutRight> @orangeB@orange("B",@redN0@red,@darkgreenN1@darkgreen);
parallel_for("TestLayoutRight",
  MDRangePolicy<Rank<2,Iterate::Right,Iterate::Right>>(
	{@red0@red,@darkgreen0@darkgreen},{@redN0@red,@darkgreenN1@darkgreen}),
  KOKKOS_LAMBDA(int @redi@red, int @darkgreenj@darkgreen) {
    @orangeB@orange(@redi@red,@darkgreenj@darkgreen) = 1000.0 * @redi@red + 1.0*@darkgreenj@darkgreen;
});
\end{code}

\pause

	\begin{block}{Default Patterns Match}
Default iteration patterns match the default memory layouts!
	\end{block}
\end{frame}



%==========================================================================

\begin{frame}[fragile]{Exercise - mdrange: Initialize multi-dim views with MDRangePolicy}

  \textbf{Details}:
  \begin{small}
  \begin{itemize}
    \item Location: \ExerciseDirectory{mdrange/Begin}
    \item This begins with the \texttt{Solution} of 02
    \item Initialize the device Views x and y directly on the device using a \texttt{parallel\_for} and \texttt{RangePolicy}
    \item Initialize the device View matrix A directly on the device using a \texttt{parallel\_for} and \texttt{MDRangePolicy}
  \end{itemize}
  \end{small}

\begin{code}
  # Compile for CPU
  cmake -B build_openmp -DKokkos_ENABLE_OPENMP=ON
  cmake --build build_openmp
  # Run on CPU
  ./build_openmp/mdrange_exercise -S 26
  # Note the warnings, set appropriate environment variables
  # Compile for GPU
  cmake -B build_cuda -DKokkos_ENABLE_CUDA=ON
  cmake --build build_cuda
  # Run on GPU
  ./build_cuda/mdrange_exercise -S 26
\end{code}

\end{frame}

%==========================================================================

\begin{frame}[fragile]{Common Policy Arguments}

  \textbf{Template Parameters common to ALL policies.}

  \begin{itemize}
     \item \texttt{ExecutionSpace}: control where code executes
     \begin{itemize}
       \item{\small{\textbf{Options:} Serial, OpenMP, Threads, Cuda, HIP, ...}}
     \end{itemize}
     \item \texttt{Schedule$<$Options$>$}: set scheduling policy.
     \begin{itemize}
        \item{\small{\textbf{Options:} Static, Dynamic}}
     \end{itemize}

     \item \texttt{IndexType$<$Options$>$}: control internal indexing type
     \begin{itemize}
       \item{\small{\textbf{Options:} int, long, etc}}
     \end{itemize}

     \item \texttt{WorkTag}: enables multiple operators in one functor
	     \begin{code}[keywords={struct,Tag1,Tag2,void,int}]
struct Foo {
  struct Tag1{}; struct Tag2{};
  KOKKOS_FUNCTION void operator(Tag1, int i) const {...}
  KOKKOS_FUNCTION void operator(Tag2, int i) const {...}
  void run_both(int N) { 
    parallel_for(RangePolicy<Tag1>(0,N),*this);
    parallel_for(RangePolicy<Tag2>(0,N),*this);
  }
});
  \end{code}
  \end{itemize}

\end{frame}

%==========================================================================

\begin{frame}[fragile]{MDRangePolicy Section Summary}

  \textbf{MDRangePolicy}
  \begin{itemize}
    \item{allows for tightly nested loops similar to OpenMP's collapse clause.}
    \item{requires functors/lambdas with as many parameters as its rank is.}
    \item{works with \texttt{parallel\_for} and \texttt{parallel\_reduce}.}
    \item{uses a tiling strategy for the iteration space.}
    \item{provides compile time control over iteration patterns.}
  \end{itemize}

\end{frame}

%==========================================================================
